\chapter*{Notation}
\addcontentsline{toc}{chapter}{Notation}
\markboth{Notation}{Notation}

The notation used throughout this book is collected here for reference.
The reader who encounters an unfamiliar symbol in the text is
encouraged to return to this section. Notation that is specific to a
single chapter is defined in that chapter and is not necessarily
listed here.

\section*{Sets and standard mathematical objects}

\begin{notationbox}[Sets and operators]
\begin{description}[leftmargin=3cm,style=nextline]
    \item[{$\mathbb{N}$}] the natural numbers, $\{0, 1, 2, \dots\}$.
    \item[{$\mathbb{Z}$}] the integers.
    \item[{$\mathbb{Q}$}] the rational numbers.
    \item[{$\mathbb{R}$}] the real numbers.
    \item[{$\mathbb{R}_+$}] the non-negative reals, $[0, \infty)$.
    \item[{$\mathbb{R}_{++}$}] the strictly positive reals, $(0, \infty)$.
    \item[{$\mathbb{C}$}] the complex numbers.
    \item[{$\mathbb{R}^n$}] the $n$-dimensional real Euclidean space.
    \item[{$\subseteq, \subset$}] subset, proper subset.
    \item[{$\cup, \cap, \setminus$}] union, intersection, set difference.
    \item[{$\emptyset$}] the empty set.
    \item[{$|S|$}] cardinality of set $S$ (number of elements).
    \item[{$f: X \to Y$}] a function from $X$ to $Y$.
    \item[{$f \circ g$}] composition of functions, $(f \circ g)(x) = f(g(x))$.
    \item[{$\lim_{x\to a}$}] limit as $x$ approaches $a$.
    \item[{$\sum, \prod, \int$}] standard sum, product, integral.
    \item[{$\partial / \partial x$}] partial derivative.
    \item[{$\nabla$}] gradient operator.
    \item[{$\propto$}] ``proportional to.''
    \item[{$\sim$}] ``scales as'' or ``has the same order as.''
    \item[{$\ll, \gg$}] much less than, much greater than.
\end{description}
\end{notationbox}

\section*{Probability and statistics}

\begin{notationbox}[Probability notation]
\begin{description}[leftmargin=3cm,style=nextline]
    \item[{$(\Omega, \mathcal{F}, P)$}] a probability space, with sample
        space $\Omega$, sigma-algebra $\mathcal{F}$, and probability
        measure $P$.
    \item[{$P(A)$}] probability of event $A \in \mathcal{F}$.
    \item[{$P(A \mid B)$}] conditional probability of $A$ given $B$.
    \item[{$X, Y, Z$}] random variables (capital Latin letters).
    \item[{$F_X(x) = P(X \leq x)$}] cumulative distribution function of $X$.
    \item[{$f_X(x)$}] probability density function of $X$ (when it exists).
    \item[{$\mathbb{E}[X]$}] expected value (mean) of $X$.
    \item[{$\text{Var}(X)$}] variance of $X$.
    \item[{$\sigma_X$ or $\text{SD}(X)$}] standard deviation of $X$.
    \item[{$\text{Cov}(X, Y)$}] covariance.
    \item[{$\text{Corr}(X, Y)$}] correlation.
    \item[{$X \sim D$}] $X$ is distributed according to $D$.
    \item[{$X_1, X_2, \dots \overset{\text{i.i.d.}}{\sim} D$}]
        independent and identically distributed.
    \item[{$\langle X \rangle_t$}] time average of $X$ over time interval $t$.
    \item[{$\langle X \rangle_n$}] ensemble average of $X$ over $n$ replicas.
    \item[{$N(\mu, \sigma^2)$}] normal distribution with mean $\mu$ and
        variance $\sigma^2$.
\end{description}
\end{notationbox}

The distinction between the time average $\langle X \rangle_t$ and
the ensemble average $\langle X \rangle_n$ is fundamental to this work
and is developed at length in Chapter 6 (in the present volume's
sequel) and applied throughout Part III. The reader should note now
that these are not in general equal, contrary to the assumption built
into much of standard economics.

\section*{Quantum mechanical notation}

We use the standard Dirac notation throughout for quantum-mechanical
discussions in Chapters 5 and 16.

\begin{notationbox}[Quantum notation]
\begin{description}[leftmargin=3cm,style=nextline]
    \item[{$|\psi\rangle$}] a quantum state vector (``ket'').
    \item[{$\langle \psi |$}] the dual vector (``bra'').
    \item[{$\langle \psi | \phi \rangle$}] inner product of quantum states.
    \item[{$|\psi|^2 = \langle \psi | \psi \rangle$}] squared norm.
    \item[{$\hat{A}, \hat{H}$}] operators (e.g., Hamiltonian).
    \item[{$\hat{A} | \psi \rangle$}] action of operator on a state.
    \item[{$\langle \hat{A} \rangle = \langle \psi | \hat{A} | \psi \rangle$}]
        expectation value of observable $\hat{A}$ in state $|\psi\rangle$.
    \item[{$\hbar$}] reduced Planck constant, $h/(2\pi) \approx 1.055 \times 10^{-34}$~J$\cdot$s.
    \item[{$|0\rangle, |1\rangle$}] computational basis states.
    \item[{$\otimes$}] tensor product.
\end{description}
\end{notationbox}

\section*{Economic notation}

We follow the conventions of Mas-Colell, Whinston, and Green (1995)
where applicable.

\begin{notationbox}[Economic notation]
\begin{description}[leftmargin=3cm,style=nextline]
    \item[{$\mathcal{X}$}] consumption set (the set of consumption bundles
        available in principle).
    \item[{$x = (x_1, \dots, x_L) \in \mathbb{R}^L_+$}] a consumption bundle
        of $L$ goods.
    \item[{$\succeq, \succ, \sim$}] preference relations: weak preference,
        strict preference, indifference.
    \item[{$u: \mathcal{X} \to \mathbb{R}$}] a utility function.
    \item[{$p = (p_1, \dots, p_L) \in \mathbb{R}^L_{++}$}] a price vector.
    \item[{$w$}] wealth (sometimes endowment); occasionally written
        $W$ when household wealth is meant.
    \item[{$B(p, w) = \{x \in \mathcal{X} : p \cdot x \leq w\}$}]
        Walrasian budget set.
    \item[{$x(p, w)$}] Walrasian demand correspondence.
    \item[{$h(p, u)$}] Hicksian (compensated) demand.
    \item[{$v(p, w)$}] indirect utility.
    \item[{$e(p, u)$}] expenditure function.
    \item[{$\beta \in (0, 1)$}] discount factor.
    \item[{$\delta$}] discount rate, often $\delta = 1 - \beta$ or
        $\beta = e^{-\delta}$ depending on context.
    \item[{$r$}] interest rate.
    \item[{$g$}] growth rate.
\end{description}
\end{notationbox}

\section*{Notation specific to this book}

Several notations are introduced for the analysis of Islamic economic
phenomena and recur throughout. They are listed here for reference and
defined in detail at first use.

\begin{notationbox}[Project-specific notation]
\begin{description}[leftmargin=3cm,style=nextline]
    \item[{$R_i(t)$}] the instantaneous \arb{rizq} (provision) of
        agent $i$ at time $t$. Used in the formal treatment of the
        rizq-conservation hypothesis.
    \item[{$K_i$}] the lifetime integral of $R_i(t)$ for agent $i$;
        the total provision allotment.
    \item[{$\mathbb{I}(H)$}] the conjunction of Islamic structural
        conditions on a household $H$, defined formally in Chapter 11.
    \item[{$\beta$}] the scaling exponent in the wealth-size relation
        $W(n) \sim n^\beta$. (Note: this is unrelated to the
        discount factor of the same symbol; context disambiguates.)
    \item[{$B_i(t)$}] the \arb{barakah} multiplier on agent $i$'s
        wealth at time $t$; defined operationally in Chapter 15 as
        the welfare residual after material controls.
    \item[{$N(t), P(t)$}] nominal and productive wealth, respectively;
        used in the riba-as-financial-entropy treatment of Chapter 13.
    \item[{$G(t) = N(t) - P(t)$}] the riba gap; the divergence of
        nominal from real wealth.
    \item[{$\nu_i$}] the \arb{niyyah} variable; an indicator of
        intentional state used in the operationalization of barakah
        and certain transaction-level analyses.
    \item[{$\mathbb{W}_t$}] absolute being, \arb{wujud}, in
        Akbarian-formal contexts. We use a distinct symbol because the
        concept does not reduce to ``existence'' in the modern sense.
    \item[{$\mathcal{T}_n$}] the $n$-th \arb{tajalli}; the
        $n$-th moment of divine self-disclosure, in the formal
        treatment of continuous creation in Chapter 9.
\end{description}
\end{notationbox}

\section*{A note on the symbol $\beta$}

The symbol $\beta$ is used in this book in two unrelated senses. In
the macroeconomic and intertemporal-choice contexts, $\beta$ denotes
the discount factor (the standard meaning). In Chapter 11 and the
related discussions of household scaling, $\beta$ denotes the
exponent in the scaling relation $W(n) \sim n^\beta$. Context will
always make clear which is intended; where the two might be confused,
explicit subscripts are used. We have not found a way to resolve this
without departing from established notation in one of the two
literatures, and have judged that established notation is the lesser
evil.

\section*{A note on the use of \arb{italic}}

Arabic terms appear in italic at first use and at points where the
term is being thematized. Once a term is established, it appears in
roman type. Mathematical variables appear in italic by the standard
convention of \LaTeX. The reader who encounters an italicized non-Latin
word should consult the glossary if unfamiliar.

\cleardoublepage
